\section{Exact multiple densities and uniform cofactor calibration}
\label{sec:tpc27-multiple-calibration}

We now estimate the factors \(B_{k,S}(m)\) in
\eqref{eq:tpc27-Ramanujan-square}.  The main cancellation is between an
exact finite-model identity and the main term of a proved uniform
M\"obius calibration theorem.

\subsection{An exact finite-model multiple identity}

Let \(m=\ell d_m\) be one retained row and put
\begin{equation}
 M_m=d_mH.
 \label{eq:tpc27-Mm}
\end{equation}
The integer \(M_m\) is squarefree, while \(\ell>R\).

\begin{proposition}[Exact model density on multiples]
\label{prop:tpc27-model-multiple}
Let \(k\) be squarefree and \((k,mH)=1\).  Then
\begin{equation}
 \boxed{
 \sum_{\substack{u\le R\\k\mid u\\(u,mH)=1}}
 \frac{\lambda'_R(u)}u
 =\frac{\mu(k)}{\varphi(k)}
 \sum_{\substack{a\mid M_m\\ak\le R}}
 \frac1{\varphi(a)}.}
 \label{eq:tpc27-truncated-model-multiple}
\end{equation}
In particular, if \(kM_m\le R\), then
\begin{equation}
 \boxed{
 \sum_{\substack{u\le R\\k\mid u\\(u,mH)=1}}
 \frac{\lambda'_R(u)}u
 =\frac{\mu(k)}{\varphi(k)}
 \rho_{M,H,R}(m).}
 \label{eq:tpc27-model-multiple}
\end{equation}
\end{proposition}

\begin{proof}
From \eqref{eq:tpc27-lambda-prime}, the left side of
\eqref{eq:tpc27-truncated-model-multiple} is
\begin{equation}
 \sum_{\substack{ub\le R\\k\mid u\\(u,mH)=1}}
 \frac{\mu(ub)\mu(b)}{\varphi(ub)}.
 \label{eq:tpc27-model-double-sum}
\end{equation}
Only squarefree \(q=ub\) contributes.  Grouping by \(q\) gives
\begin{equation}
 \sum_{q\le R}\frac{\mu(q)}{\varphi(q)}
 \sum_{\substack{b\mid q\\k\mid q/b\\(q/b,mH)=1}}
 \mu(b).
 \label{eq:tpc27-model-group-q}
\end{equation}
Because \(\ell>R\), no \(q\le R\) contains the large row prime.  Put
\(a=(q,M_m)\).  The condition \((q/b,M_m)=1\) forces \(a\mid b\).
Since \((k,M_m)=1\), write \(q=a q_0\) and \(b=a b_0\).  The inner
sum becomes
\begin{equation}
 \mu(a)\sum_{b_0\mid q_0/k}\mu(b_0),
 \label{eq:tpc27-inner-Mobius-collapse}
\end{equation}
and is zero unless \(q_0=k\), in which case it equals \(\mu(a)\).
Thus \(q=ak\) with \(a\mid M_m\), and the surviving summand is
\begin{equation}
 \frac{\mu(ak)\mu(a)}{\varphi(ak)}
 =\frac{\mu(k)}{\varphi(k)\varphi(a)}.
\end{equation}
This proves \eqref{eq:tpc27-truncated-model-multiple}.  If
\(kM_m\le R\), every divisor \(a\mid M_m\) occurs, and
\begin{equation}
 \sum_{a\mid M_m}\frac1{\varphi(a)}
 =\prod_{p\mid M_m}\left(1+\frac1{p-1}\right)
 =\frac{M_m}{\varphi(M_m)}
 =\rho_{M,H,R}(m).
\end{equation}
This proves \eqref{eq:tpc27-model-multiple}.
\end{proof}

The cutoff condition in this proposition is why the opened-row bound
\(d_m\le R^{1/2}\) matters.  It leaves a fixed-power range of
multipliers \(k\) on which the model identity is complete.

\subsection{The matching prime multiple}

Define
\begin{equation}
 A_{k,S}^{\rm mult}(m)
 :=\sum_{\substack{u\le S\\k\mid u\\(u,mH)=1}}
 \frac{a(u)}u.
 \label{eq:tpc27-prime-multiple-def}
\end{equation}

\begin{lemma}[Prime multiple calibration]
\label{lem:tpc27-prime-multiple}
Let \(k\) be squarefree and \((k,mH)=1\).  Then the exact identity
\begin{equation}
 A_{k,S}^{\rm mult}(m)
 =\frac{\mu(k)}k
 \left\{
 A_{kmH}(S/k)-(\log k)C_{kmH}(S/k)
 \right\}
 \label{eq:tpc27-prime-multiple-exact}
\end{equation}
holds.  Fix \(0<\kappa<s\), where
\(S=X^{s+o(1)}\).  Uniformly for \(k\le X^\kappa\) and every fixed
\(A>0\),
\begin{equation}
 \boxed{
 A_{k,S}^{\rm mult}(m)
 =\frac{\mu(k)}{\varphi(k)}\rho_{P,H}(m)
 +O_{A,\mathscr D}\!\left(
 \frac{\rho_{P,H}(m)(1+\log k)}
 {\varphi(k)(\log X)^A}
 \right).}
 \label{eq:tpc27-prime-multiple}
\end{equation}
\end{lemma}

\begin{proof}
On the squarefree support, write \(u=kw\).  Then
\((w,kmH)=1\) and
\begin{equation}
 a(kw)=\mu(k)a(w)-\mu(k)\mu(w)\log k.
 \label{eq:tpc27-a-kw}
\end{equation}
Division by \(kw\) and summation prove
\eqref{eq:tpc27-prime-multiple-exact}.

For \(k\le X^\kappa\), the calibration length
\(Y=S/k\) is at least a fixed positive power of \(X\).  Also
\(kmH\le X^{O(1)}\), so one fixed constant \(C\), depending only on
the exponent data, gives \(kmH\le Y^C\).  Apply
\cref{eq:tpc27-uniform-A,eq:tpc27-uniform-C} with a calibration
exponent larger than the requested \(A\).  Since \((k,mH)=1\), the
main term is
\begin{equation}
 \frac{\mu(k)}k\frac{kmH}{\varphi(kmH)}
 =\frac{\mu(k)}{\varphi(k)}\rho_{P,H}(m).
\end{equation}
The two error terms are bounded by the displayed remainder because
\(\log(S/k)\asymp\log X\) uniformly in this fixed-power range.  This
proves \eqref{eq:tpc27-prime-multiple}.
\end{proof}

\subsection{Small and large multiple densities}

Set
\begin{equation}
 K_0=X^{\kappa+o(1)},
 \qquad
 0<\kappa<\min\left\{s,\frac{r_0}{2}\right\}.
 \label{eq:tpc27-K0}
\end{equation}
Because \(d_m\le R^{1/2}\) and \(H\) is fixed, all sufficiently large
\(X\) satisfy
\begin{equation}
 kM_m\le R
 \qquad(k\le K_0).
 \label{eq:tpc27-small-k-model-range}
\end{equation}

\begin{corollary}[Calibrated small multiples]
\label{cor:tpc27-small-Bk}
For every retained row, every effective squarefree \(k\le K_0\), and
every fixed \(A>0\),
\begin{equation}
 \boxed{
 B_{k,S}(m)
 =\frac{\mu(k)}{\varphi(k)}\delta_H(m)
 +O_{A,\mathscr D}\!\left(
 \frac1{\varphi(k)(\log X)^A}
 \right).}
 \label{eq:tpc27-small-Bk}
\end{equation}
\end{corollary}

\begin{proof}
By \eqref{eq:tpc27-bR}, \(B_{k,S}(m)\) is the prime multiple
\eqref{eq:tpc27-prime-multiple-def} minus the model multiple.  Apply
\cref{prop:tpc27-model-multiple,lem:tpc27-prime-multiple} and use
\eqref{eq:tpc27-row-drift}.  Since
\(\rho_{P,H}(m)\ll_h\log\log(3X)\) and
\(1+\log k\ll\log X\), invoke
\cref{lem:tpc27-prime-multiple} with three additional logarithmic
powers and then rename the exponent.  This absorbs the explicit
polylogarithmic factors without introducing an \(X^\eps\) loss.
\end{proof}

For large \(k\), cancellation is unnecessary.  We import the elementary
cofactor mass lemma of TPC-24: uniformly for squarefree \(k\le S\),
\begin{equation}
 \sum_{\substack{w\le S/k\\(w,k)=1}}
 \frac{|b_R(kw)|}{w}
 \ll_{C_S}(\log(2X))^4.
 \label{eq:tpc27-cofactor-L1}
\end{equation}
It follows from the positive divisor formula for \(\lambda'_R\) and
harmonic summation, and uses no cancellation
\citep{WangTPC24}.

\begin{lemma}[Absolute large-multiple bound]
\label{lem:tpc27-large-Bk}
Uniformly for every row and every \(k\ge1\),
\begin{equation}
 \boxed{
 |B_{k,S}(m)|\ll_{\mathscr D}
 \frac{(\log(2X))^4}{k}.}
 \label{eq:tpc27-large-Bk}
\end{equation}
\end{lemma}

\begin{proof}
There is nothing to prove unless \(k\) is squarefree and
\((k,mH)=1\).  Write \(u=kw\) in \eqref{eq:tpc27-Bks}.  Squarefree
support forces \((w,k)=1\), while the remaining row-coprimality
condition only decreases the absolute sum.  Thus
\begin{equation*}
 |B_{k,S}(m)|
 \le\frac1k
 \sum_{\substack{w\le S/k\\(w,k)=1}}
 \frac{|b_R(kw)|}{w},
\end{equation*}
and \eqref{eq:tpc27-cofactor-L1} proves the result.
\end{proof}

The two bounds \eqref{eq:tpc27-small-Bk} and
\eqref{eq:tpc27-large-Bk} concern the same exact one-row factor.  The
split at \(K_0\) therefore loses no intermediate term.
